\documentclass[11pt]{article}

\usepackage{amsmath,amssymb,amsthm}
\usepackage{enumitem}
\usepackage[margin=1in]{geometry}
\usepackage{setspace}

\onehalfspacing

% Numbering setup
\theoremstyle{definition}
\newtheorem{definition}{Definition}[section]

\theoremstyle{plain}
\newtheorem{axiom}{Axiom}[subsection]
\newtheorem{theorem}{Theorem}[subsection]

\setlist[itemize]{topsep=0.5em,itemsep=0.25em}

\begin{document}

\section*{Introduction}

These notes present several axiomatic systems for geometry. We begin with abstract
\emph{practice geometries}, designed to emphasize how different axiom sets generate
different mathematical worlds.

Certain foundational terms are intentionally left undefined. This reflects the axiomatic
method: objects are characterized not by intrinsic meaning, but by the relations imposed
by the axioms.

\section{Practice Geometries}

The following terms will be left undefined:
\begin{itemize}
    \item point
    \item line
    \item lie on
    \item passing through
\end{itemize}

\subsection*{Definitions}

\begin{definition}
Points $A$ and $B$ are \emph{collinear} if there exists a line that passes through both $A$ and $B$.
\end{definition}

\begin{definition}
Lines $\ell$ and $m$ \emph{intersect} if there exists a point $A$ such that $A$ lies on both $\ell$
and $m$.
\end{definition}

\subsection{Practice Geometry \#1}

\subsubsection*{Axioms}

\begin{axiom}
There exist four points such that no three are collinear.
\end{axiom}

\begin{axiom}
For each pair of points, there exists a unique line passing through both points.
\end{axiom}

\begin{axiom}
Every pair of lines intersect.
\end{axiom}

\begin{axiom}
There exists a line with exactly five points on it.
\end{axiom}

\subsubsection*{Theorems}

\begin{theorem}
Every pair of lines intersect at exactly one point.
\end{theorem}

\begin{theorem}
There exist four lines such that no three intersect at the same point.
\end{theorem}

\begin{theorem}
There exists a point with exactly five lines passing through it.
\end{theorem}

\begin{theorem}
For any pair of points, there exists a line that does not pass through either point.
\end{theorem}

\begin{theorem}
Each point lies on exactly five lines.
\end{theorem}

\begin{theorem}
Each line passes through exactly five points.
\end{theorem}

\begin{theorem}
There are exactly 21 points.
\end{theorem}

\begin{theorem}
There are exactly 21 lines.
\end{theorem}

\subsection{Practice Geometry \#2}

\subsubsection*{Axioms}

\begin{axiom}
There exist four points such that no three are collinear.
\end{axiom}

\begin{axiom}
For each pair of points, there exists a unique line passing through both points.
\end{axiom}

\begin{axiom}
Given a line $\ell$ and a point $P$ not on $\ell$, there is a unique line through $P$
that does not intersect $\ell$.
\end{axiom}

\begin{axiom}
There exists a line with exactly three points on it.
\end{axiom}

\subsubsection*{Theorems}

\begin{theorem}
Every point lies on exactly four lines.
\end{theorem}

\begin{theorem}
Each line passes through exactly three points.
\end{theorem}

\begin{theorem}
For each line $\ell$, there are exactly two lines that do not intersect $\ell$.
\end{theorem}

\begin{theorem}
There are exactly nine points.
\end{theorem}

\begin{theorem}
There are exactly twelve lines.
\end{theorem}


\end{document}